\section*{Bab \ref{ch:g6:latitude-longitude} --- Lintang, Bujur, dan Zona Waktu}

\begin{solution}{exo:g6:latitude-longitude:1}
``\omterm{def:g6:latitude-longitude:latitude}{Lintang} diukur dari \omterm{def:g3:continents-and-oceans:equator}{khatulistiwa} dan bernilai dari $0^\circ$ sampai
$90^\circ$ (utara atau selatan); \omterm{def:g6:latitude-longitude:longitude}{bujur} diukur dari \omterm{def:g6:latitude-longitude:longitude}{meridian utama} dan
bernilai dari $0^\circ$ sampai $180^\circ$ (timur atau barat).''
\end{solution}

\begin{solution}{exo:g6:latitude-longitude:2}
\omterm{def:g3:continents-and-oceans:poles}{Kutub Utara}: $90^\circ$\,LU. \omterm{def:g3:continents-and-oceans:poles}{Kutub Selatan}: $90^\circ$\,LS. Kota di
\omterm{def:g3:continents-and-oceans:equator}{khatulistiwa}: $0^\circ$.
\end{solution}

\begin{solution}{exo:g6:latitude-longitude:3}
Sydney: \omterm{def:g3:continents-and-oceans:equator}{belahan bumi selatan}, timur. New York: \omterm{def:g3:continents-and-oceans:equator}{belahan bumi utara}, barat.
Kairo: \omterm{def:g3:continents-and-oceans:equator}{belahan bumi utara}, timur. Rio de Janeiro: \omterm{def:g3:continents-and-oceans:equator}{belahan bumi selatan},
barat.
\end{solution}

\begin{solution}{exo:g6:latitude-longitude:4}
(a) Benar. (b) Salah: \omterm{def:g6:latitude-longitude:latitude}{paralel} makin pendek ke arah \omterm{def:g3:continents-and-oceans:poles}{kutub}; yang paling
panjang adalah \omterm{def:g3:continents-and-oceans:equator}{khatulistiwa}. (c) Benar. (d) Salah: \omterm{def:g6:latitude-longitude:longitude}{meridian utama} melalui
Greenwich, di London.
\end{solution}

\begin{solution}{exo:g6:latitude-longitude:5}
Kairo. New York. Tokyo terletak pada kira-kira $36^\circ$\,LU,
$140^\circ$\,BT (sedikit di atas \omterm{def:g6:latitude-longitude:latitude}{paralel} $30^\circ$\,LU, sedikit melewati
\omterm{def:g6:latitude-longitude:longitude}{meridian} $120^\circ$\,BT, belum sampai $150^\circ$\,BT).
\end{solution}

\begin{solution}{exo:g6:latitude-longitude:6}
$45 \div 15 = 3$: tengah hari kedua kota itu berselisih 3~jam. Tengah hari
terjadi lebih dahulu di kota~B, yang letaknya lebih ke timur.
\end{solution}

\begin{solution}{exo:g6:latitude-longitude:7}
$73 \div 15 \approx 4.9$, dibulatkan menjadi 5: zonanya $+5$. Ketika di
Greenwich pukul 12.00 siang, di Mumbai pukul $12 + 5 = 17$, yaitu pukul
17.00.
\end{solution}

\begin{solution}{exo:g6:latitude-longitude:8}
$30 \times 111 = 3\,330$~km. \omterm{def:g3:continents-and-oceans:poles}{Kutub Utara} terletak pada $90^\circ$:
$90 \times 111 = 9\,990$~km, kira-kira 10\,000~km.
\end{solution}

\begin{solution}{exo:g6:latitude-longitude:9}
Selisihnya $9 - (-5) = 14$~jam. $8 + 14 = 22$: di Tokyo pukul 22.00 hari
Senin.
\end{solution}

\begin{solution}{exo:g6:latitude-longitude:10}
Ia benar. Bumi berputar dari barat ke timur, jadi kota-kota di timur lebih
dahulu terbawa menghadap Matahari: kota-kota itu melihat Matahari terbit,
mencapai titik tertingginya, dan terbenam lebih dahulu daripada kota-kota
yang lebih ke barat.
\end{solution}

\begin{solution}{exo:g6:latitude-longitude:11}
(a) $42 \times 3 = 126$~detik, sedikit lebih dari 2~menit. (b)
$126 \div 240 \approx 0.5$: kira-kira setengah derajat \omterm{def:g6:latitude-longitude:longitude}{bujur}. (c)
$0.5 \times 111 \approx 55$~km (dengan $0.525$: kira-kira 58~km). Kesalahan
sekitar lima puluh kilometer: cukup untuk tidak menemukan sebuah pulau,
atau untuk menabrak \omterm{def:g1:landscapes-around-us:water}{pesisir} yang dikira masih jauh.
\end{solution}

\begin{solution}{exo:g6:latitude-longitude:12}
(a) $18 - 9 = 9$: pukul 09.00 hari Senin di Greenwich. (b) $9 + 10 = 19$:
pukul 19.00 hari Senin di Greenwich. (c) $19 - 8 = 11$: pukul 11.00 hari
Senin di Los Angeles. Pesawat itu berangkat pada pukul 18.00 hari Senin
(waktu Tokyo) dan mendarat pada pukul 11.00 hari Senin yang sama (waktu Los
Angeles): jam di tempat tiba menunjukkan waktu yang lebih awal daripada jam
di tempat berangkat, karena pesawat itu menyeberangi \omterm{def:g6:latitude-longitude:date-line}{garis batas tanggal} ke
arah timur dan mundur satu hari.
\end{solution}

\begin{solution}{pb:g6:latitude-longitude:1}
\textbf{1.} Ana (Rio de Janeiro, $23^\circ$\,LS) dan Mia (Auckland,
$37^\circ$\,LS).

\textbf{2.} Ana dan Sam ($43^\circ$\,BB dan $118^\circ$\,BB). Lena berada di
\omterm{def:g6:latitude-longitude:longitude}{meridian utama}.

\textbf{3.} Lena: London terletak pada $52^\circ$\,LU, lebih jauh dari
\omterm{def:g3:continents-and-oceans:equator}{khatulistiwa} daripada semua kota lainnya.

\textbf{4.} Mia: Auckland terletak pada $175^\circ$\,BT, hanya $180 - 175 = 5$
derajat dari \omterm{def:g6:latitude-longitude:longitude}{meridian} $180^\circ$.

\textbf{5.} $360^\circ \div 24 = 15^\circ$.

\textbf{6.} London: $0$. Rio de Janeiro: $43 \div 15 \approx 2.9$, barat,
jadi $-3$. Tokyo: $140 \div 15 \approx 9.3$, timur, jadi $+9$. Auckland:
$175 \div 15 \approx 11.7$, timur, jadi $+12$. Los Angeles:
$118 \div 15 \approx 7.9$, barat, jadi $-8$.

\textbf{7.} Rio de Janeiro: $20 - 3 = 17$, pukul 17.00 hari Jumat. Tokyo:
$20 + 9 = 29$; $29 - 24 = 5$: pukul 05.00 hari Sabtu.

\textbf{8.} Tokyo terletak jauh di timur: jamnya 9~jam lebih dahulu daripada
jam London. Pada pukul 20.00 di London, Tokyo sudah melewati tengah malam,
jadi hari Sabtunya sudah dimulai.

\textbf{9.} $20 - 8 = 12$: pukul 12.00 siang hari Jumat.

\textbf{10.} $12 - (-8) = 20$~jam lebih dahulu.

\textbf{11.} Ketika menyeberangi \omterm{def:g6:latitude-longitude:date-line}{garis batas tanggal} ke arah timur, kapal itu
mundur ke hari sebelumnya.

\textbf{12.} $20 + 12 = 32$; $32 - 24 = 8$: ketika panggilan dimulai, di
Auckland \textbf{pukul 08.00 hari Sabtu}. Mia bergabung dalam panggilan pada
waktu sarapan, sementara Sam makan siang pada hari Jumat.
\end{solution}
